\documentclass[14pt, a4paper]{extreport}

\usepackage{kpi-lab}
% Defined by subject
\newcommand{\CourseTitle}{Проектування програмних систем для мобільних пристроїв}
\newcommand{\Teacher}{асистент, Нестерук А. О.}
\newcommand{\ReportYear}{2026}
% Defined by Student
\newcommand{\StudentName}{Ковальов Олександр}
\newcommand{\StudentGroup}{ІМ-51мн}
\newcommand{\Variant}{4}
\newcommand{\CourseNumber}{2}
% Defined by Submission
\newcommand{\Topic}{Емулятор та ресурси Android}
\newcommand{\LabNumber}{4}

\begin{document}
	\pagenumbering{gobble}
	
	\begin{titlepage}
		\begin{center}
			\normalsize
			Національний технічний університет України\\
			«Київський політехнічний інститут імені Ігоря Сікорського» \\[0.5em]
			Факультет інформатики та обчислювальної техніки\\
			Кафедра обчислювальної техніки
			
			\vfill
			
			{\large \textbf{ЗВІТ}\\[0.5em]}
			{з лабораторної роботи №\LabNumber \\}
			{з дисципліни <<\CourseTitle>>}
			
			\vspace{2em}
			
			{\textbf{<<\Topic>>}}
			
			\vspace{1em}
			
			{\textbf{Варіант №\Variant}}
			
			\vfill
			
			\begin{flushright}
				Виконав: \\ студент \CourseNumber{} курсу, групи \StudentGroup \\
				\StudentName \\[1em]
				Перевірив: \\ \Teacher
			\end{flushright}
			
			\vfill
			
			\begin{flushright}
				Дата здачі: \DTMtoday
			\end{flushright}
			
			\vspace{5em}
			КИЇВ -- \ReportYear
		\end{center}
	\end{titlepage}
	
	\clearpage
	\pagenumbering{arabic}
	\setcounter{page}{2}
	
	\subsection*{Мета} 
	
	Вивчити програмні засоби емуляції операційної системи Android, розглянути її ресурси, основні методи роботи із файловою системою, ознайомитись з методами підключення бібліотек та віджетів у додаток, вивчити маніфест операційної системи.
	
	\subsection*{Завдання}
	
	Додати елементи управління (\texttt{Button}, \texttt{TextView}, \texttt{EditText}, \texttt{ListView}) в програму, розроблену в лабораторній роботі №3. Запустити програму на емуляторі.
	
	\subsection*{Завдання за варіантом}
	
	Доменна область: Абітурієнт. 
	
	Поля:  ID, Прізвище, Ім'я, По-батькові, Адреса, Телефон, Оцінки. 
	
	Створити масив об'єктів. 
	
	Вивести:
	
	\begin{enumerate}[label=\alph*.]
		\item список абітурієнтів, які мають незадовільні оцінки;
		
		\item список абітурієнтів, середній бал у яких вище заданого;
		
		\item вибрати задане число \texttt{N} абітурієнтів, що мають найбільш високій середній бал (вивести також повний список абітурієнтів, які мають напівпрохідний бал).
	\end{enumerate}
	
	\section*{Хід роботи.}

	Для виконання роботи використано застосунок, розроблений у третій лабораторній роботі, без зміни архітектури MVC чи стеку -- Rust, Tauri, React з TypeScript. Для даної роботи створено новий окремий екран -- <<Пошук>>, на якому всі чотири елементи (\texttt{Button}, \texttt{TextView}, \texttt{EditText} та \texttt{ListView}) використано явно.
	
	\image{9cm}{home-screen}
	
	На боці Rust додано контролер-модуль \texttt{search} із відповідною командою, яка фільтрує список абітурієнтів репозиторію за входженням підрядка в прізвище без врахування регістру і повертає результат у тому ж вигляді, що й список абітурієнтів на екрані <<Абітурієнти>>. На боці React створено відповідний <<екран>>, що поєднує всі чотири елементи.
	
	\texttt{EditText} -- текстове поле, контрольоване React-станом через атрибути \texttt{value} та \texttt{onChange}. Це функціональний аналог \texttt{TextWatcher}/\texttt{addTextChangedListener} в Android -- кожна зміна тексту одразу оновлює стан JS, так само як \texttt{TextWatcher} сповіщає Activity про зміну вмісту поля.
	
	\texttt{Button} -- кнопка <<Пошук>>, що за натисканням викликає відповідну функцію. Це аналог \texttt{setOnClickListener} на \texttt{android.widget.Button}.
	
	\texttt{ListView} -- результати виведено не таблицею, як на інших екранах, а списком клікабельних елементів, по одному на абітурієнта. Клік на елементі зберігає його \texttt{ID} у стані \texttt{selectedId}. Це ближчий аналог \texttt{android.widget.ListView} з \texttt{Adapter}'ом та \texttt{setOnItemClickListener}, ніж таблиця -- один елемент інтерфейсу на один рядок даних, з власним обробником кліку.
	
	\texttt{TextView} -- під списком розміщено окрему секцію <<Деталі>>, яка виводить повну інформацію про вибраного абітурієнта (ID, ПІБ, адреса, телефон, оцінки, середній бал) лише для читання. Блок оновлюється декларативно при зміні \texttt{selectedId}, що є React-аналогом імперативного виклику \texttt{textView.setText(...)} всередині обробника кліку на елементі списку в Android.
	
	\image{10cm}{page}
	
	Коректність реалізації перевірено на емуляторі Android -- уведено частину прізвища в поле пошуку, натиснуто кнопку <<Пошук>>, і список скоротився до абітурієнтів, чиє прізвище містить введений підрядок.
	
	\image{9cm}{search}
	
	Після натискання на один з елементів списку під ним з'являється детальна інформація про вибраного абітурієнта.
	
	\image{9cm}{details}

	\subsection*{Висновок}

	У цій лабораторній роботі до застосунку <<Абітурієнт>> додано елементи керування \texttt{Button}, \texttt{TextView}, \texttt{EditText} та \texttt{ListView}, явно виокремлені в новий екран <<Пошук>>. Наведено відповідність кожного React-елемента нативному Android-віджету. Коректність роботи підтверджено на емуляторі Android.

	\section*{Контрольні запитання.}

	\begin{enumerate}
		\item \textbf{Дайте визначення поняття <<емулятор>>.} Це апаратна або програмна система, яка точно моделює поведінку іншої системи, дозволяючи виконувати програми або використовувати периферійні пристрої, призначені для оригінальної архітектури, на цільовій платформі. Головна мета емуляції полягає у відтворенні внутрішньої структури та логіки роботи пристрою на низькому рівні, що відрізняє її від симуляції, яка лише імітує зовнішню абстрактну поведінку.
		
		\item \textbf{Назвіть основні етапи встановлення емулятора Android.} Процес встановлення зазвичай відбувається через інтегроване середовище Android Studio, де спочатку необхідно відкрити SDK Manager для завантаження потрібних версій Android SDK та системних образів, а потім перейти до Device Manager для створення нового віртуального пристрою, обравши апаратний профіль, версію операційної системи та налаштувавши параметри продуктивності, такі як обсяг оперативної пам'яті та графічне прискорення.
		
		\item \textbf{Наведіть види емуляторів. Назвіть їх переваги та недоліки.} Емулятори поділяються на апаратні та програмні, а також за рівнем системи на емулятори архітектури процесора та емулятори повноцінної операційної системи. Їхня головна перевага полягає у можливості тестувати програмне забезпечення на різноманітних конфігураціях обладнання без необхідності купувати реальні фізичні пристрої, що суттєво економить кошти та спрощує процес зневадження, проте недоліком є висока вимогливість до ресурсів хостової машини, можливе зниження швидкодії через трансляцію інструкцій та іноді неповна або неточна підтримка специфічних апаратних функцій, таких як датчики чи камери.
		
		\item \textbf{Назвіть основні ресурси Android.} Основні ресурси в Android-додатках зберігаються у каталозі \texttt{res} і включають макети користувацького інтерфейсу у форматі XML, строкові змінні для інтернаціоналізації, графічні зображення різної роздільної здатності, кольори, стилі та теми, а також анімації і сирі бінарні файли, що в сукупності дозволяє суворо відокремити логіку програмного коду від представлення та полегшує адаптацію програми під різні конфігурації екранів пристроїв.
		
		\item \textbf{Назвіть методи роботи з файлами.} Робота з файлами в Android базується на стандартних класах мови Java або Kotlin з використанням потоків вводу-виводу, при цьому операційна система надає специфічні методи контексту, наприклад \texttt{openFileInput} та \texttt{openFileOutput} для безпечного доступу до внутрішнього ізольованого сховища додатку, а також окремі інструменти та API для роботи із зовнішнім сховищем, кеш-пам'яттю та файлами з каталогу ресурсів або \texttt{assets}.
		
		\item \textbf{Охарактеризуйте формат зберігання даних JSON.} Це легкий текстовий формат обміну даними, який легко читається людьми та ефективно парситься машинами, що базується на двох основних структурах: колекціях пар ключ-значення, які представляють об'єкти, та впорядкованих структурах значень, які представляють масиви, що робить його універсальним стандартом для передачі даних у веб-сервісах та клієнт-серверній архітектурі незалежно від мови програмування.
		
		\item \textbf{Назвіть основні етапи підключення бібліотек до Android Studio.} Підключення зовнішніх бібліотек найчастіше здійснюється через систему збирання Gradle, що вимагає відкриття файлу \texttt{build.gradle} на рівні модуля, додавання необхідної залежності з вказанням її координат у блок \texttt{dependencies} та подальшої синхронізації проекту для завантаження артефактів, хоча також можливе ручне додавання локальних компільованих файлів безпосередньо в директорію \texttt{libs} із відповідним оновленням конфігурації збирання.
		
		\item \textbf{Назвіть призначення файлу \texttt{manifest.xml} та охарактеризуйте його.} Файл \texttt{AndroidManifest.xml} є ключовим конфігураційним документом будь-якого Android-додатку, який надає операційній системі життєво важливу метаінформацію перед запуском коду, визначаючи простір імен програми, її головні компоненти, такі як активності, фонові сервіси та приймачі широкомовних повідомлень, а також чітко декларуючи необхідні дозволи для доступу до захищених API, апаратних функцій та мінімальні вимоги до версії SDK.
	\end{enumerate}

	\pagebreak

	\section*{Лістинг.}

	\textbf{search.rs}

	\begin{minted}{rust}
		//! Lab 4: a separate screen that explicitly
		//! demonstrates `Button`, `TextView`, `EditText`, and `ListView`.
		
		// State<T> must be taken by value in every command here — that's Tauri's
		// calling convention, not a real ownership need.
		#![allow(clippy::needless_pass_by_value, reason = "Tauri's convention")]
		
		use crate::applicant::ApplicantView;
		use crate::repository::Repository;
		use std::sync::{Mutex, PoisonError};
		use tauri::State;
		
		#[tauri::command]
		pub fn search_applicants(
		state: State<'_, Mutex<Repository>>, query: String,
		) -> Vec<ApplicantView> {
			let repo = state.lock().unwrap_or_else(PoisonError::into_inner);
			let needle = query.trim().to_lowercase();
			
			repo.list()
			.iter()
			.filter(|a| needle.is_empty() || a.last_name.to_lowercase().contains(&needle))
			.map(ApplicantView::from)
			.collect()
		}
	\end{minted}
	
	\textbf{Search.tsx}
	
	\begin{minted}{tsx}
		import { useEffect, useState } from "react";
		import { invoke } from "@tauri-apps/api/core";
		import { Link } from "react-router-dom";
		import type { Applicant } from "../types";
		
		// Demonstrates a Button, TextView, EditText, and ListView on a single screen (Lab 4 assignment).
		// An EditText (search field) and a Button ("Search") filter the ListView (the list below);
		// clicking an item in the list populates
		// the TextView (a read-only details section) — all four elements are visible simultaneously.
		export function Search() {
			const [query, setQuery] = useState("");
			const [rows, setRows] = useState<Applicant[] | null>(null);
			const [selectedId, setSelectedId] = useState<number | null>(null);
			
			async function search() {
				const result = await invoke<Applicant[]>("search_applicants", { query });
				setRows(result);
				if (!result.some((a) => a.id === selectedId)) {
					setSelectedId(null);
				}
			}
			
			useEffect(() => {
				search();
			}, []);
			
			const selected = rows?.find((a) => a.id === selectedId) ?? null;
			
			return (
			<main className="container">
			<Link className="back" to="/">
			← На головну
			</Link>
			<h1>Пошук абітурієнтів</h1>
			<p className="subtitle">
			EditText + Button + ListView + TextView в одному екрані
			</p>
			
			{/* EditText + Button */}
			<form
			className="controls"
			onSubmit={(e) => {
					e.preventDefault();
					search();
			}}
			>
			<label>
			Прізвище містить
			<input
			type="text"
			value={query}
			onChange={(e) => setQuery(e.currentTarget.value)}
			/>
			</label>
			<button type="submit">Пошук</button>
			</form>
			
			{/* ListView */}
			<section className="panel">
			<h2>
			Результати <span className="count">({rows?.length ?? 0})</span>
			</h2>
			{rows && rows.length > 0 ? (
				<ul className="list-view">
				{rows.map((a) => (
					<li
					key={a.id}
					className={a.id === selectedId ? "selected" : undefined}
					onClick={() => setSelectedId(a.id)}
					>
					{a.last_name} {a.first_name} {a.patronymic} — сер. бал{" "}
					{a.average}
					</li>
					))}
				</ul>
				) : (
				<p className="empty">— немає —</p>
				)}
			</section>
			
			{/* TextView */}
			<section className="panel">
			<h2>Деталі</h2>
			{selected ? (
				<p className="text-view">
				ID: {selected.id}
				<br />
				ПІБ: {selected.last_name} {selected.first_name}{" "}
				{selected.patronymic}
				<br />
				Адреса: {selected.address}
				<br />
				Телефон: {selected.phone}
				<br />
				Оцінки: {selected.grades.join(", ")}
				<br />
				Середній бал: {selected.average}
				</p>
				) : (
				<p className="empty">Виберіть абітурієнта зі списку</p>
				)}
			</section>
			</main>
			);
		}
	\end{minted}
\end{document}